\documentclass[../../main.tex]{subfiles}

\begin{document}

\section{赋范线性空间}\label{section:ch13-赋范线性空间}

线性空间 $X$ 是一个阿贝尔群. 群的加法运算记为 $+$, 给定实数 $\alpha$ 和 $u\in X$, 可定义数乘 $\alpha\cdot u\in X$ 使得以下三个性质成立: 对于实数 $\alpha$ 和 $\beta$ 以及 $X$ 中的成员 $u$ 和 $v$,
\begin{align*}
(\alpha+\beta)\cdot u=\alpha\cdot u+\beta\cdot u,\\
\alpha\cdot(u+v)=\alpha\cdot u+\alpha\cdot v,\\
(\alpha\beta)\cdot u=\alpha\cdot(\beta\cdot u)\text{ 且 }1\cdot u=u.
\end{align*}
线性空间也称为向量空间, 遵循 $\mathbf{R}^{n}$, 线性空间的成员常称为向量. 线性空间的典型例子是任意非空集 $D$ 上的实值函数族, 对两个函数 $f,g:D\to\mathbf{R}$ 与实数 $\lambda$, 加法 $f+g$ 与数乘 $\lambda\cdot f$ 在 $D$ 上逐点定义为
\begin{align*}
\text{对所有 }x\in D,\quad(f+g)(x)=f(x)+g(x)\text{ 而 }(\lambda\cdot f)(x)=\lambda f(x).
\end{align*}
回忆我们首先在第 7 章研究过的线性空间 $X$ 上的范数概念, 定义在线性空间 $X$ 上的非负实值函数 $\|\cdot\|$ 称为范数, 若对所有 $u,v\in X$ 与 $\alpha\in\mathbf{R}$,
\begin{align*}
\|u\|=0\text{ 当且仅当 }u=0,\\
\|u+v\|\leqslant\|u\|+\|v\|,\\
\|\alpha u\|=|\alpha|\|u\|.
\end{align*}
如同我们在第 9 章观察到的, 线性空间上的范数诱导了该空间上的度量, 其中 $u$ 与 $v$ 之间的距离定义为 $\|u-v\|$. 当我们谈到赋范空间的度量性质, 如有界性与完备性, 我们指的是由范数诱导的度量. 类似地, 当谈到拓扑性质, 如序列收敛或集合的开、闭或紧, 我们指的是由上述度量诱导的拓扑.

\begin{definition}[等价范数]\label{definition:ch13-等价范数}
线性空间 $X$ 上的两个范数 $\|\cdot\|_{1}$ 与 $\|\cdot\|_{2}$ 称为等价的, 若存在常数 $c_{1}\geqslant 0$ 和 $c_{2}\geqslant 0$ 使得
\begin{align*}
\text{对所有 }x\in X,\quad c_{1}\cdot\|x\|_{1}\leqslant\|x\|_{2}\leqslant c_{2}\cdot\|x\|_{1}.
\end{align*}
\end{definition}

我们立即看到两个范数等价当且仅当它们诱导的度量等价. 因此, 若线性空间上的一个范数被等价的范数代替, 拓扑和度量性质仍然不变.

有限维空间的代数的概念对一般的线性空间也很重要. 给定线性空间 $X$ 中的向量 $x_{1},\cdots,x_{n}$ 与实数 $\lambda_{1},\cdots,\lambda_{n}$, 向量
\begin{align*}
x=\sum_{k=1}^{n}\lambda_{k}x_{k}
\end{align*}
称为 $x_{i}$ 的线性组合. $X$ 的非空子集 $Y$ 称为线性子空间或简单地称为子空间, 若 $Y$ 中的每个线性组合也属于 $Y$.

对于 $X$ 的非空子集 $S$, 谈到 $S$ 的扩张我们指的是 $S$ 中的向量的所有线性组合: 我们将 $S$ 的扩张记为 $\operatorname{span}[S]$. 我们把证明 $\operatorname{span}[S]$ 是 $X$ 的线性子空间留作习题. 它在以下意义下是包含 $S$ 的最小线性子空间: 它包含于任何包含 $S$ 的线性子空间. 若 $Y=\operatorname{span}[S]$, 我们说 $S$ 张成 $Y$. 考虑 $S$ 的扩张的闭包 $\overline{\operatorname{span}[S]}$ 是有用的. 我们把 $X$ 的线性子空间的闭包是线性子空间留作练习. 因此 $\overline{\operatorname{span}[S]}$ 是 $X$ 的线性子空间. 它在以下意义下是包含 $S$ 的最小闭线性子空间: 它包含于任何包含 $S$ 的闭线性子空间. 我们称 $\overline{\operatorname{span}[S]}$ 为 $S$ 的闭线性扩张.

对于线性空间 $X$ 的任何两个非空子集 $A$ 与 $B$, 我们定义 $A$ 与 $B$ 的和 $A+B$ 为:
\begin{align*}
A+B=\{x+y\mid x\in A,\ y\in B\}.
\end{align*}
在 $B$ 是单点集 $\{x_{0}\}$ 的情形中, 我们记 $A+\{x_{0}\}$ 为 $A+x_{0}$. 称该集为 $A$ 的一个平移. 对 $\lambda\in\mathbf{R}$, 我们定义 $\lambda A$ 为形如 $\lambda x$(其中 $x\in A$)的所有元素的集合. 观察到若 $Y$ 和 $Z$ 是 $X$ 的子空间, 则和 $Y+Z$ 也是 $X$ 的子空间. 在 $Y\cap Z=\{0\}$ 的情形中, 我们记 $Y+Z$ 为 $Y\oplus Z$, 且称 $X$ 的这个子空间为 $Y$ 与 $Z$ 的直和.

对于赋范线性空间 $X$, 中心在原点、半径为 $1$ 的开球 $\{x\in X\mid\|x\|<1\}$ 称为 $X$ 中的单位开球, 而 $\{x\in X\mid\|x\|\leqslant 1\}$ 称为 $X$ 中的单位闭球. 我们称满足 $\|x\|=1$ 的向量 $x\in X$ 为单位向量.

几乎所有关于度量空间的重要定理都要求完备性. 因此不足为奇的是, 在赋范线性空间中那些关于由范数诱导的度量是完备的空间最为重要.

\begin{definition}[Banach 空间]\label{definition:ch13-Banach}
一个赋范线性空间称为 Banach 空间, 若它作为具有由范数诱导的度量的度量空间是完备的.
\end{definition}

Riesz-Fischer 定理告诉我们对于实数的可测集 $E$ 和 $1\leqslant p\leqslant\infty$, $L^{p}(E)$ 是一个 Banach 空间. 我们也证明了对于紧拓扑空间 $X$, 具有最大值范数的 $C(X)$ 是一个 Banach 空间. 当然, 我们从关于 $\mathbf{R}$ 的完备性公理推出每个欧氏空间 $\mathbf{R}^{n}$ 是一个 Banach 空间.

\begin{exercise}\label{exercise:ch13-01}
证明线性空间 $X$ 的非空子集 $S$ 是一个子空间当且仅当 $S+S=S$, 且对每个 $\lambda\in\mathbf{R}$, $\lambda\neq 0$, $\lambda\cdot S=S$.
\end{exercise}

\begin{exercise}\label{exercise:ch13-02}
若 $Y$ 与 $Z$ 是线性空间 $X$ 的子空间, 证明 $Y+Z$ 也是子空间且 $Y+Z=\operatorname{span}[Y\cup Z]$.
\end{exercise}

\begin{exercise}\label{exercise:ch13-03}
令 $S$ 为赋范线性空间 $X$ 的子集.
\begin{enumerate}[(i)]
\item 证明 $X$ 的线性子空间族的交也是 $X$ 的线性子空间.
\item 证明 $\operatorname{span}[S]$ 是 $X$ 的所有包含 $S$ 的线性子空间的交, 因此是 $X$ 的一个线性子空间.
\item 证明 $\overline{\operatorname{span}[S]}$ 是 $X$ 的所有包含 $S$ 的闭线性子空间的交, 因此是 $X$ 的一个闭线性子空间.
\end{enumerate}
\end{exercise}

\begin{exercise}\label{exercise:ch13-04}
对赋范线性空间 $X$, 证明函数 $\|\cdot\|:X\to\mathbf{R}$ 是连续的.
\end{exercise}

\begin{exercise}\label{exercise:ch13-05}
对两个赋范线性空间 $(X,\|\cdot\|_{1})$ 与 $(Y,\|\cdot\|_{2})$, 通过令 $\lambda\cdot(x,y)=(\lambda x,\lambda y)$ 与 $(x_{1},y_{1})+(x_{2},y_{2})=(x_{1}+x_{2},y_{1}+y_{2})$ 在笛卡儿积 $X\times Y$ 上定义线性结构. 对 $x\in X$ 与 $y\in Y$ 定义乘积范数 $\|\cdot\|$ 为 $\|(x,y)\|=\|x\|_{1}+\|y\|_{2}$. 证明这是一个范数, 序列关于该范数收敛当且仅当两个分量序列都收敛. 进一步地, 证明: 若 $X$ 和 $Y$ 都是 Banach 空间, 则 $X\times Y$ 也是.
\end{exercise}

\begin{exercise}\label{exercise:ch13-06}
令 $X$ 为赋范线性空间.
\begin{enumerate}[(i)]
\item 令 $\{x_{n}\}$ 与 $\{y_{n}\}$ 为 $X$ 中的序列使得 $\{x_{n}\}\to x$ 且 $\{y_{n}\}\to y$. 证明对任何实数 $\alpha$ 与 $\beta$, $\{\alpha x_{n}+\beta y_{n}\}\to\alpha x+\beta y$.
\item 用(i)证明: 若 $Y$ 是 $X$ 的子空间, 则它的闭包 $\overline{Y}$ 也是 $X$ 的线性子空间.
\item 用(i)证明: 向量的和是从 $X\times X$ 到 $X$ 的连续映射而数乘是从 $\mathbf{R}\times X$ 到 $X$ 的连续映射.
\end{enumerate}
\end{exercise}

\begin{exercise}\label{exercise:ch13-07}
证明 $[a,b]$ 上的所有多项式组成的集合 $\mathcal{P}$ 是线性空间. 把 $\mathcal{P}$ 考虑为赋范线性空间 $C[a,b]$ 的子集, 证明 $\mathcal{P}$ 不是闭的. 把 $\mathcal{P}$ 考虑为赋范线性空间 $L^{1}[a,b]$ 的子集, 证明 $\mathcal{P}$ 不是闭的.
\end{exercise}

\begin{exercise}\label{exercise:ch13-08}
定义在向量空间 $X$ 上的非负实值函数 $\|\cdot\|$ 称为拟范数, 若 $\|x+y\|\leqslant\|x\|+\|y\|$ 且 $\|\alpha x\|=|\alpha|\|x\|$. 若 $\|x-y\|=0$, 定义 $x\cong y$. 证明这是一个等价关系. 定义 $X/\cong$ 为 $X$ 的在 $\cong$ 下的等价类的集合且对 $x\in X$ 定义 $[x]$ 为 $x$ 的等价类. 若我们定义 $\alpha[x]+\beta[y]$ 为 $\alpha x+\beta y$ 的等价类且定义 $\|[x]\|=\|x\|$, 证明 $X/\cong$ 是赋范向量空间. 用 $X=L^{p}[a,b](1\leqslant p<\infty)$ 说明这一程序.
\end{exercise}

\end{document}